\documentclass[12pt]{article}

\usepackage{amssymb}

\title{Well-ordered choice and the weak partition principle fail in New Foundations}

\author{Elliott Glazer and Lavinia Randall Holmes\footnote{The second author is the same person as the Randall Holmes, M. Randall Holmes, or M.R. Holmes mentioned in the bibliography}}

\begin{document}

\maketitle

\begin{abstract}

Specker showed in 1954 (\cite{speckernotac}) that the Axiom of Choice is false in Quine's set theory New Foundations (introduced in the paper \cite{nf}).  It has been unclear ever since how or whether this result could be extended.  This paper reports substantial progress:  the axiom of well-ordered choices (each well-ordered partition has a choice set) and the weak partition principle (and so of course the partition principle proper) are inconsistent with NF.  The paper also contains a proof of Infinity in NF which does not use ordered pairs, demonstrating that the type level ordered pair commonly used in NF, whose properties depend on Infinity holding, can be shown to work without preliminary use of the usual Kuratowski pair definition:  this fixes an ugliness in the usual exposition of basic concepts in NF.

\end{abstract}

\subsection{Version notes}

\begin{description}

\item[10/24/2026:]  This version contains a careful re-edit of the proof of Henson's Lemma.

\item[10/23/2026:]  This version now contains a short proof of Infinity followed by a justification of the claim that this is a proof of the same kind as the proofs appearing earlier, or can be construed as such.

\item[10/15/2026:]  This version defers considerations about Infinity and the original Specker proof to the end (these sections are not prepared yet).  The aim is to write the proof of Infinity in a way which conforms to
the structure of the new proofs, and to compare the original Specker proof with the new proofs in such a way as to bring out similarities and differences.  A final section discussing Specker trees is proposed.

\item[10/12/2026:]  Moved the proof of Infinity to the end, and recast it to make the analogy with the proofs of the new choice refutations clear.  Moved Henson's lemma to appear first.  I project a further section on Specker's original results and a section on Specler trees.

\item[10/8/2026:]  added a discussion of the Partition Principle, the Weak Partition Principle, and the consequence of WPP which we explicitly refute.

\item[10/7/2026:]  added a remark as to why the proof of Infinity doesnt work in TST, and minor edits.

\item[10/7/2026:]  This copy is intended to replace strict inequality with non-strict inequality in the Henson section, with correlated effects on inequalities used.

Also, the version notes are in the right place :-)

Also, some inequalities in the Henson argument seem to have been reversed:  I believe I have corrected them.

\item[10/7/2026:]  Inserted an ordered-pair-free proof of Infinity, which seems relevant to the project at hand [one needs Infinity to justify the Quine pair which rationalizes the exponents on T's in the proof].
Also, though this is not brought out in the current text explicitly, the closing moves of the proof of Infinity are another example of Henson's maneuver.  It might be that Henson's maneuver should appear earlier,
or the relevance to the proof of Infinity could be brought out retrospectively.  I do not think a proof of Infinity which does not use pairing has appeared before, though I would have to check.  If it is thought too cumbersome, note that the role of section 2.1 in this draft is entirely as an extended comment on the paragraph immediately before the section.

I added more commentary to clarify the place of section 2.1 in the development, and an indication of why the proof in section 2.1 is of the same kind as the arguments in the later section, in the Henson section.  I think I was to rewrite the Henson section using $F(T(\alpha)) \geq \alpha$ rather than $F(T(\alpha)) > \alpha$;  I believe this is straightforward, and the maneuvers I use to make the inequality strict are artificial, and obscure matters a little in the comparison with the proof of Infinity.

\item[10/6/2026, 1:14 pm Boise time:]  The first release of a draft.  This contains pointers to citations but no citations:  these are to be filled in later.  Do we want to add a further discussion of Specker trees, or
discussion of the original Glazer inequality?  Mostly this says, I am actually working on it :-)



\end{description}

\subsection{Statement about the use of artificial intelligence in the background of this paper, and about the contributions of the human authors}

No part of this paper was written using AI, nor are any of the main results of the paper due to artificial intelligence.  However, the convenient form of the disproof of the Partition Principle found in section 4 of this paper was found by ChatGPT Sol.  This result was originally proved using a different inequality found by the first author whose proof involved elaborate forcing technology.  The second author was convinced by the simple form of the statement of the original inequality that it ought to have an elementary proof not involving forcing techniques.  The first author asked ChatGPT Sol to find such a proof.   ChatGPT Sol did not find an elementary proof
of the first author's identity.  Instead, it formulated a different inequality which would serve the same purpose, proved it by elementary methods, and pointed out that it would serve the purpose.  The proof as presented in this paper was written by hand by the second author, but the role of AI here is of considerable interest.

The first result of the paper (the disproof of well-ordered choice) was completely carried out by the first author, though not in the form presented here.  The first author proved that the weak partition principle fails in NF plus the Axiom of Counting (the assertion that the set of natural numbers is strongly cantorian, formulated by Rosser in \cite{rosserbook}).  The second author adapted methods of C Ward Henson (\cite{hensontheorem}) ultimately derived from contemplation of Specker's disproof of the Axiom of Choice (\cite{speckernotac}) to simplify the disproof of well-ordered choice and complete the disproof of the weak partition principle in NF.   The technical difficulty of this paper was greatly reduced by ChatGPT Sol's observation that the first author's inequality used to disprove the weak partition principle could be replaced by the one used here, as noted above.  The second author provided the ordered-pair-free proof of Infinity in NF.

\section{The theory of types and New Foundations}

Quine's set theory New Foundations (NF), so called after the title of the paper \cite{nf} in which it was proposed in 1937, is most efficiently defined and motivated by first defining an extremely streamlined typed theory of sets which we call TST.  The history of TST is confusing:  a description of TST is often incorrectly given as a description of the far more complicated type theory of the {\em Principia Mathematica\/} of Russell and Whitehead.  Norbert Wiener seems to have first observed the possibility of a theory like TST in \cite{wienerpair}, 1914:  the crucial simplification of Russell's theory was the possibility of defining an ordered pair in terms of sets.  An actual formal description of TST (looking very different from this one) seems to have been given first by Tarski in \cite{tarskitst} in 1930, only seven years before the definition of New Foundations.

The modern definition of TST follows.  It is a sorted first order theory with equality and membership as primitive relations.  The sorts are indexed by the natural numbers.   We do not require that our variables be adorned with type superscripts:  we denote the type of a variable $x$ by ${\tt type}(x)$.  The condition for well-formedness of a formula $x=y$ is that ${\tt type}(x)={\tt type}(y)$.  The condition for well-formedness of a formula $x \in y$
is that ${\tt type}(x)+1={\tt type}(y)$.  The axiom schemes of the theory are extensionality [every formula of the shape $(\forall xy:x=y \leftrightarrow (\forall z:z \in x \leftrightarrow z \in y))$ is an axiom] and comprehension
[every formula of the shape $(\exists A:(\forall x:x \in A \leftrightarrow \phi))$, where $A$ does not appear in $\phi$, is an axiom:  the unique witness to this formula may be denoted by $\{x\mid \phi\}$.].  An interesting thing to note is that these are exactly the axioms of naive set theory, with the inconsistencies averted by the restrictions on formation of formulas imposed by the typing scheme.

What Quine noticed in 1937 (following observations of Russell about his much more complicated type theory) is that TST has a symmetry property which suggests a simplification of the theory.
If we postulate an injective map on variables , whose restriction to variables of each type $i$ is a bijection onto variables of type $i+1$, and denote the image of a variable $x$ under this map by $x^+$, we can conveniently define $\phi^+$ for each formula $\phi$ in the language of TST as the result of replacing each variable $u$ appearing in $\phi$ with $u^+$.  This clearly preserves well-formedness.  Quine's observation can be summed up as pointing out that, for any formula $\phi$, if we can prove $\phi$, we can also prove $\phi^+$ [the converse is not true] and any object we can defined as a set abstract $\{x\mid \phi\}$ has a precisely analogous object $\{x^+\mid \phi^+\}$ in the next higher type.  Quine made the leap of suggesting that we can in fact simply identify the types.

The theory NF is a single-sorted theory with equality and membership as primitive relations.   Its axioms are $(\forall xy:x=y \leftrightarrow (\forall z:z \in x \leftrightarrow z \in y))$ [a single extensionality axiom]
and every formula of the shape $$(\exists A:(\forall x:x \in A \leftrightarrow \phi)),$$ where $A$ does not appear in $\phi$, and where there is a function $\sigma$ from variables to natural numbers such that
in each atomic subformula $u=v$ of $\phi$ we have $\sigma(u)=\sigma(v)$ and in each atomic subformula $u \in v$ of $\phi$ we have $\sigma(u)+1=\sigma(v)$.  A formula $\phi$ for which there is such a function $\sigma$ is said to be {\em stratified\/} with stratification $\sigma$, and this axiom scheme is called {\em stratified comprehension\/}.  The intended effect here is that the axioms of NF are exactly the axioms stated in single sorted
language which are the result of dropping type distinctions between variables in an axiom of TST without introducing any identifications of variables.   The axiom of stratified comprehension is often dismissed as
a ``syntactical trick''.  This impression can be dispelled by observing that stratified comprehension is finitely axiomatizable, so the theory can be presented without even indirect reference to the language of TST.

We note [contra those who think there is some technical difficulty with term constructions in NF] that while the language of NF, like that of ZF, does not contain any term constructions as primitive, term constructions are readily handled, and moreover, readily handled relative to the stratification criterion.  A discussion of this is found in the original paper.  Basically, the theory of definite descriptions handles this efficiently.
We define $(\theta x:\phi)$ as the unique object $x$ such that $\phi$, if there is one, and otherwise as the unique object [of the same type as $x$, in the TST context] $\emptyset$ with no elements.  In NF or TST, this works straightforwardly:  any occurrence
of $(\theta x:\phi)$ can be eliminated as if defined contextually, following Russell (the refinement of providing a default value when the description would otherwise be undefined is due to Rosser).  In TST, $(\theta x:\phi)$ is assigned the same type as $x$, and the typing rules for atomic formulas generalize in the obvious way.
In NF, $(\theta x:\phi)$ is to be assigned the same value as $x$ by a stratification, and if $\phi$ is stratified, contextual expansion of $(\theta x:\phi)$ is easily seen to preserve stratification.  We define $\{x\mid \phi\}$
as $(\theta A:(\forall x:x \in A \leftrightarrow \phi))$, where $A$ does not occur in $\phi$, and that is all that needs to be said about term constructions.

Specker showed in \cite{speckeramb}, 1962, that Quine's leap of faith to identification of the types is at least partially justified.  A more cautious conjecture to make on the basis of the symmetry of TST described above
is to adopt the axiom scheme $\phi \leftrightarrow \phi^+$ for each closed formula $\phi$, which is called the Ambiguity Scheme.  Specker showed that TST + Ambiguity is equiconsistent with NF.

But Specker had already showed, in \cite{speckernotac}, 1954, that the Axiom of Choice is false in NF.  This was surprising and alarming, and had the corollary that NF is significantly stronger than TST.  It is provable in TST that every finite set is well-orderable, so of course this is provable in NF, so from $\neg$AC we can conclude that the universe is not well-orderable [if it were, every partition of the universe would obviously have a choice set] and so infinite.   TST does not prove Infinity, so NF is essentially stronger.

We will not need to give Specker's disproof of Choice at the outset, as we can prove the stronger result in a subsequent section that the axiom of well-ordered choice is false in NF.

It has proved remarkably difficult to extend Specker's result or conversely to determine what weakenings of Choice might be consistent with NF.  This paper is the first significant extension of Specker's original results of 1954 in a negative direction.


\section{Preliminaries on the mathematics of relations, functions, and cardinality in TST and NF}

In this section, we will give a brief description of how mathematics is done in TST and NF, focussing in particular on definitions which will be unfamiliar to those working in ZF or ZFC.

We begin by defining the natural numbers.  The basic idea (following Frege) is that the natural number $n$ is the collection of all sets with $n$ elements, which appears circular.  We demonstrate that it is not.

We define $\emptyset$ as $\{x\mid \neg x=x\}$.  We define $V$ as $\{x\mid x=x\}$.

We define $\iota(x)$ or $\{x\}$ as $\{y\mid y=x\}$.

We define $A \cup B$ as $\{x\mid x \in A \vee x \in B\}$.

For any term $T(x_1,\ldots,x_n)$ containing just those free variables, we define $\{T(x_1,\ldots,x_n\mid \phi\}$ as $\{x\mid (\exists x_1,\ldots,x_n:x=T(x_1,\ldots,x_n) \wedge \phi\}$.

For any term construction $F(x)$ with one free variable, define $F``A$ as $\{F(a)\mid a \in A\}$.

For any set $A$, we define $\bigcup A$ as $\{x\mid (\exists a \in A:x \in a)\}$ and $\bigcap A$ as $\{x:(\forall a\in A:x \in a)\}$.

We define 0 as $\{\emptyset\}$.  Notice that this is the collection of all sets with zero elements.

For any set $A$, we define $\sigma(A)$ as $\{a \cup \{x\}:a \in A \wedge x \not\in a\}$.  Notice that if $A$ were the set of all sets with $n$ elements, $\sigma(A)$ would be the set of all sets with $n+1$ elements.

We define $\cal I$ as $\{A:0 \in A \wedge (\forall x:x \in A \rightarrow \sigma(x)\in A)\}$.  This is the collection of inductive sets.

We define $\mathbb N$ as $\{n:(\forall I \in {\cal I}:n \in I\}$, the intersection $\bigcap \cal I$ of all inductive sets.

These definitions can be construed in two different ways.  In NF, they work exactly as written (note that all of them are stratified).  In TST, we must be aware that each of these is a definition scheme which defines objects in each sufficiently high type.
$\emptyset$ for example is defined in each type $i+1$.  0 is definable in each type $i+2$.  $\sigma(A)$ is defined and of type $i+2$ for each $A$ of type $i+2$.  $\cal I$ is defined in each type $i+4$ and
$\mathbb N$ is defined in each type $i+3$.  In general, we will give definitions in the more convenient single sorted way allowed by NF, but all of our definitions will be stratified and  usable with care in TST.

We define $\mathbb F$, the collection of finite sets, as $\bigcup \mathbb N$.

We state the Axiom of Infinity, which must be adjoined to TST as an additional axiom scheme, but which is actually a difficult theorem of NF, as $V \not\in \cal F$.  In what follows immediately, we simply assume this.  In a subsequent section, we will prove it, in a way which brings out the similarity of the proof of Infinity to the proofs of the refutations of choice which we present first.

We then use the natural numbers to define the ordered pair.  The usual definition of $(x,y)$ as $\{\{x\},\{x,y\}\}$ can be used in TST (with or without Infinity) or NF, but is formally inconvenient because $(x,y)$ is then two types higher than $x$ or $y$. 

Instead, we give the following definitions (following Quine in \cite{quinepair}:  the use of this pair has been customary in NF since it was proposed) .  Define $\sigma^*(x)$ as $\sigma(x)$ if $x \in \mathbb N$ and as $x$ otherwise.  It is a theorem that $\sigma^*$ is an injection.  Note that $\sigma$ itself is not an injection on sets generally (we leave this as an exercise).   Define $\sigma_1(A)$ as $\sigma^*``A$ and $\sigma_2(A)$ as $(\sigma^*``A) \cup \{0\}$.  Then define $(A,B)$ as $\sigma_1``A \cup \sigma_2``B$.  This is an ordered pair,
because $\sigma_1$ and $\sigma_2$ are injective with disjoint ranges, and $(A,B)$ is of the same type as $A$ and $B$.  It is further entertaining though we will not make any use of it that every set is an ordered pair under this definition.  We may in our subsequent development comment in passing on ways the development would differ (in being somewhat more inconvenient) if the Kuratowski pair were used.  [Of course, in our subsequent proof of Infinity, we will {\em not\/} use this ordered pair definition:  in fact, we will contrive to use no ordered pair definition at all.]

Once we are in possession of the ordered pair, we can define basic concepts of the theories of relations and functions just as usual.  It is particularly worthy of note that $A \times B = \{(a,b):a \in A \wedge b \in B\}$ is a set
(which it would be if we used the Kuratowski pair as well) but further is a set of the same type as the common type of $A$ and $B$, which would not be the case if the Kuratowski pair were used, which makes the definitions
of cardinal addition and multiplication much easier to state.

We define the cardinality $|A|$ of a set $A$ as $\{B\mid A \sim B\}$, where $A \sim B$ is defined as usual as asserting the existence of a bijection from $A$ to $B$.  Notice that $|A|$ is one type higher than $A$.

An operation which is unfamiliar in ordinary set theory:  $T(|X|)$ is defined as $|\iota``X|$.  In the context of TST, this might be thought of as sending a cardinal to ``the same'' cardinal one type up.   But in NF it is a nontrivial operation, and not a set function, as we will show in the next paragraph.  It is straightforward to show that $T$ commutes with all relations and operations of cardinal arithmetic.  $T$ is injective, so the notation
$T^{-1}(|X|)$ makes sense, but we will see below that the $T^{-1}$ operation is partial.

As usual, we say that $|X| \leq |Y|$ iff there is $Z$ such that $Y \supseteq Z \sim X$.  The Schr\"oder-Bernstein theorem $|X| \leq |Y| \wedge |Y| \leq |X| \rightarrow |X| = |Y|$ admits the standard proof
in TST and so in NF.  The Cantor theorem as transliterated from ordinary set theory, $|X| < |{\cal P}(X)|$, does not make sense in TST and is false in NF (certainly $|V| < |{\cal P}(V)|$ does not hold), but this is due to a misreading.  The correct form of the Cantor theorem is $T(|X|) = |\iota``X| < |{\cal P}(X)|$, and the proof is basically standard.  In NF, we find that $T(|V|) = |\iota``V| < |{\cal P}(V)| = |V|$, from which we see that the singleton operation is not a set function and the $T$ operation is nontrivial.  $T^{-1}(|V|)$ is seen not to exist.  It is fairly direct to show further that the $T$ operation on cardinals is not a set function.  Considerations in this paragraph
motivate the definition $\exp(|X|) = T^{-1}(|{\cal P}(|X|)|)$ of the (partial) set function of exponentiation of cardinals [this definition is repeated below], and $|X| < \exp(|X|)$ is shown to hold where the right side of the inequality is defined.

We define the order type ${\tt ot}(W)$ of a well-ordering $W$ [it should be noted that we construe well-orderings as reflexive relations:  otherwise we would be unable to distinguish the ordinals 0 and 1] as $\{W'\mid W \approx W'\}$, where $W \approx W'$ asserts that $W$ and $W'$ are isomorphic binary relations (and of course implies that $W'$ is also a well-ordering).

If $W$ is a binary relation, $W^\iota$ is defined as $\{(\{x\},\{y\}):(x,y)\in W\}$.  If $W$ is a well-ordering, so is $W^\iota$, and we define $T({\tt ot}(W))$ as ${\tt ot}(W^\iota)$.  Like the $T$ operation on cardinals, the $T$ operation on ordinals amounts to projecting an ordinal upward one type.  We will immediately show that the $T$ operation on ordinals is nontrivial.

We can define ${\tt ot}(W) \leq_{\tt ord} {\tt ot}(W')$ as holding exactly if there is $W''$ such that $W' \supseteq W'' \approx W$.  This is the natural well-ordering on the ordinals (the proof that it is a well-ordering
we leave as an exercise).   Define $\Omega$ as ${\tt ot}(\leq_{\tt ord})$.

We observe that the order type of the restriction of $\leq_{\tt ord}$ to ordinals $\beta <_{\tt ord} \alpha$ is $T^2(\alpha)$.  Examination of types in the definitions reveals this.   So the order type
of the restriction of $\leq_{\tt ord}$ to the set of ordinals less than $\Omega$ is $T^2(\Omega)$, and we see that $T^2(\Omega) <_{\tt ord} \Omega$.  Did the reader think we were about to fall afoul of the Burali-Forti paradox?  This shows that the T operation on ordinals is nontrivial, and also that it is not a set function (as otherwise it would be an isomorphism between a well-ordering and a proper initial segment of that well-ordering).  It has the model theoretic consequence that the ordinals of a model of NF are never well-ordered by the natural order from the standpoint of the metatheory.   It is worth noting that the crucial statement 
$T^2(\Omega) <_{\tt ord} \Omega$, which is weird and wonderful in the NF context, makes the much more sensible assertion in TST that the sequence of ordinals in any type is shorter than the sequence of ordinals two types up.

We prove Hartog's theorem in a version appropriate TST or NF (a generalization of what we showed in the previous paragraph).  Let $X$ be a set and consider the collection $Y$ of all ${\tt ot}(W)$ for $W$ such that
${\tt dom}(W) \subseteq X$.  Let $\Omega_X$ be the order type of the natural order on the ordinals restricted to this collection.  $\Omega_X$ is two types higher than the ${\tt ot}(W)$'s in the collection.   From this we observe that $T^{-2}(\Omega(X))$ is the supremum of the ordinals in $Y$, if it exists, and greater than any ordinal in $Y$.  Further, the domain of an element of $T^{-2}(\Omega(X))$ cannot be smaller than or the same size as $X$,
as otherwise $T^{-2}(\Omega(X))$ would belong to $Y$.   The size of the domain of an element of $T^{-2}(\Omega(X))$ is $T^{-2}``|Y|$, for which we introduce the convenient notation $\aleph(|X|)$ [this definition is repeated below] and we have shown that $\aleph(|X|)$ (if defined) is not less than or equal to $|X|$.

We remind the reader that a cardinal or ordinal is referred to as {\em cantorian\/} iff it is a fixed point of the appropriate T operation.
\section{A lemma due to C W Henson}


\begin{description}
\item[Definition:]  Define a special function on cardinals of well-ordered sets as a partial function $F$ which commutes with $T$ [for all $\beta$, $F(T(\beta))=T(F(\beta))$ if $F(\beta)$ is defined], is nondecreasing (and defined on an initial segment of the well-ordered cardinals), and satisfies $F(\alpha) \geq \alpha$ for all $\alpha$ for which it is defined.

\item[Lemma:]  For no special function can we have $F(T(\alpha)) \geq \alpha$, if $T(\alpha)<\alpha$ and $\alpha$ is not dominated by a cantorian $\beta=T(\beta)$. [This result follows from Theorem 3.2 in \cite{henson}:  the fact that Henson's theorem is stated for ordinals makes no essential difference; this is merely a matter of correct attribution, as our proof is self-contained].

\item[Proof:]  We prove this.  Suppose otherwise for a particular special $F$ and cardinal $\alpha$ of a well-ordered set.

Let $\gamma$ be smallest such that for some $n$, $F^n(\gamma) \geq \alpha$.  There is such a $\gamma$ because $T(\alpha)$ is a $\delta$ such that $F^1(\delta)\geq\alpha$.

Note that $T(\gamma)$ is smallest such that for some $n$, $F^n(T(\gamma)) \geq T(\alpha)$:  if $\mu<T(\gamma)$ and $F^n(\mu) \geq T(\alpha)$ [taking $n$ smallest], then for some $\nu<\gamma$, $\mu=T(\nu)$ and for some $m$, $n=T(m)$ [because it indexes a sequence of distinct images under $T$ (except possibly the last one)], so we have $F^{T(m)}(T(\nu))\geq T(\alpha)$ and so $F^m(\nu) \geq \alpha$, which contradicts minimality of $\gamma$.   We then see that $T(\gamma)$ is the same cardinal as $\gamma$:
clearly $F^{n+1}(T(\gamma)) \geq \alpha$ so $T(\gamma)\geq \gamma$, and we have $m$ such that $F^m(\gamma) \geq \alpha > T(\alpha)$, so $T(\gamma) \leq \gamma$.   That is, $\gamma$ is cantorian.

Now let $N$ be smallest such that $F^N(\gamma)\geq \alpha$.  $N >0$ since $F^0(\gamma) = \gamma$, being cantorian, is certainly less than $\alpha$.

Certainly $T(N)$ is smallest such that $F^{T(N)}(\gamma) \geq T(\alpha)$.

We have $F^{N-1}(\gamma)<\alpha$ and $F^N(\gamma) \geq \alpha$.

We have $F^{T(N)-1}(\gamma)<T(\alpha)<\alpha$, $F^{T(N)}(\gamma) \geq T(\alpha)$, and by hypothesis $F^{T(N)+1}(\gamma) \geq \alpha$ (or $F^{T(N)+1}(\gamma)$ might be undefined, but this would imply $N=T(N)$).

But this means that $N$ must be either $T(N)$ or $T(N)+1$.

We cannot have $N=T(N)$ because this would make $F^{T(N)}(\gamma) = F^N(\gamma)$ cantorian and greater than $\alpha$.

We cannot have $N=T(N)+1$ because $N$ and $T(N)$ must have the same parity (odd or even).

We have arrived at a contradiction, completing the proof of the Lemma.



\end{description}




\section{Refuting choice from well-ordered partitions}

The axiom of well-ordered choice (AC$_{\tt wo}$) is the assertion that every well-ordered partition has a choice set.
Given an arbitrary ordinal indexed sequence of sets $X_\alpha$, we can construct a sequence of disjoint
sets $Y_{T^3\alpha}=\iota^3``X_\alpha \times \{\alpha\}$ (if $s$ is a well-ordering, we define $s_{\alpha}$ as the item in the well-ordering with segment before it of order type $\alpha$: $s_\alpha$ is of type one lower than that of the segment, which is of type one lower than the ordinal index $\alpha$:  it is useful to note that in the NF context (or the context of stratified set theory generally) an ordinal indexing is not a function).
A choice set from the well-ordered partition obtained by dropping empty sets from the range of this sequence readily gives us a well-orderable set which meets each nonempty set in the original sequence in a single element.

For any cardinal $\kappa$ and $X \in \kappa$, we define $\exp(\kappa)$ as $T^{-1}(|{\cal P}(X)|)$, so we have $T(\exp(|X|) = |{\cal P}(X)|$:  of course, this does not depend on the choice of $X \in \kappa$.  Note that this will  fail to be defined if ${\cal P}(X)$ is not at least as small as the set of singletons.   

We define $\aleph(\kappa)$, where $\kappa$ is a cardinal and $X \in \kappa$,  as $$T^{-2}(|\{{\tt ot}(W):{\tt dom}(W) \subseteq X\}|),$$ so that $T^2(\aleph(|X|)) = |\{{\tt ot}(W):{\tt dom}(W) \subseteq X\}|$:  this again does not depend on the choice of $X \in \kappa$ and may fail to be defined if the collection of order types of well-orderings of subsets of $X$  is not sufficiently small.  These are precise analogues of cardinal operations in ordinary set theory and have familiar properties mod stratification considerations.  We remind the reader that we are using a type level ordered pair:  otherwise the exponent in the definition of $\aleph$ would be different.

We also define $\aleph^*(\kappa)$ [not needed until the section after next] as $$T^{-3}(|\{{\tt ot}(W):{\tt dom}(W)\subseteq {\cal P}(X) \wedge \emptyset \not\in {\tt dom}(W) \wedge (\forall AB \in {\tt dom}(W):A \cap B = \emptyset)\})|):$$  this is the set of order types of well-orderings of partitions of any $X \in \kappa$, suitably type shifted.  Again this can be undefined if the set of order types is insufficiently small.

We compute a bound on $\aleph(\exp(\kappa))$ for any cardinal $\kappa$.   Suppose that $X \in \kappa$.
Let $Y \subseteq {\cal P}(X)$ have a well-ordering.   Then $Y \times Y$ has a well-ordering.  Let $C_\alpha = (A_\alpha,B_\alpha)$ be an ordinal indexing of $Y\times Y$.  From this derive an ordinal indexing $D$ of the set
$\{A-B:A \in Y \wedge B\in Y$ defined as $$D=\{(A-B,G-H):{\tt min}\{\alpha:A\setminus B = A_\alpha-B_\alpha\} \leq {\tt min}\{\alpha:G\setminus H=G_\alpha-H_\alpha\}\}.$$  We can use AC$_{\tt wo}$ to select a set $E$ which is a well-orderable subset of $X$ meeting each nonempty $D_\alpha$ in a single element.  The well-ordering is $F=\{(e,f):\exists AB \in {\tt dom}(D):e \in A \wedge f \in B \wedge (A,B) \in D)\}$ The map $(A \in Y \mapsto A \cap E)$ is an injection, because if $A$ and $B$ are distinct elements of $Y$, there is an element of $A \Delta B$ in $E$ by construction.  We see that any well-ordering of a set of subsets of $X$ is isomorphic to a well-ordering of subsets of a well-ordered subset $E$ of $X$:
So we have $\{{\tt ot}(W):W \subseteq {\cal P}(X)\} \subseteq \{{\tt ot}(W):W \subseteq {\cal P}(E)\}$ 
so $T^2(\aleph(T(\exp(\kappa)))) \leq T^2(\aleph(T(\exp(|E|))))$.  Now $T^2(|E|)$ is the cardinality of the set
of order types of initial segments of $F$, which is less than $T^2(\aleph(|E|))$ which is in its turn less than $T^2(\aleph(|X|))$, so we can complete the computation:  moving $T$'s about on both sides of the previous inequality, $T^2(\aleph(T(\exp(|X|)))) = \aleph(\exp(T^3(|X|))) \leq \aleph(\exp(T^3(|E|)) \leq \aleph(\exp(T^3(\aleph(|X|))))$, and from this we can see $\aleph(\exp(|X|)) \leq \aleph(\exp(\aleph(|X|)))$ (if both sides exist).

For sufficiently large concrete $n$ (large enough that all computations on both sides succeed) we have 

$\aleph(T^{n-1}(|V|))=\aleph(\exp(T^n(|V|))) \leq \aleph(\exp(\aleph(T^n(|V|))))$. 

And this is impossible, by Henson's lemma:  $F(\alpha) = \aleph(\exp(\alpha))$ is a special function on cardinals of well-ordered sets and $\alpha = \aleph(T^{n-1}(|V|))$ satisfies $T(\alpha)<\alpha$ and is not dominated by any cantorian cardinal of well-ordered sets.



So we have arrived at a contradiction.  So our original assumption of choice from well-ordered partitions is inconsistent with NF.

Note that this is a disproof of the Axiom of Choice itself, and might reveal more about what is really going on than the original proof.

It is worth observing that the novel part of this argument is precisely that any well-ordering of subsets of a set $X$ is actually a well-ordering of subsets of a well-ordered subset of $X$:  this is what was not known.



\section{Refuting the weak partition principle}

The Partition Principle (framed in a way which respects the type constraints of TST) is the assertion that for any partition $P$ of a set $X$, $|P| \leq T(|X|)$.  This is easily seen to be a consequence of the axiom of choice:  the singleton image of a choice set of $P$ will be a subset of $\iota``X$.  The Weak Partition Principle (WPP) is the assertion that for any partition $P$ of a set $X$, $$\neg |P|>T(|X|).$$  

A consequence of the Weak Partition Principle is that if $Y$ maps onto $X$ and $Y$ injects into $X$ then $X$ injects into $Y$ (i.e., $X$ and $Y$ are the same size by Schr\"oder-Bernstein).  If $Y$ maps onto $X$, there is a partition of $Y$ which is the same size as $\iota``X$.  $Y$ injects into $X$ so $\iota``Y$ injects into the partition of $Y$.  The partition of $Y$ cannot by WPP be strictly larger than $\iota``Y$ so it must be exactly the same size as $\iota``Y$, and so $X$ is the same size as $Y$.

We then argue that WPP implies that for any $X$, $$\aleph(|X \times \mathbb N|)=\aleph^*(|X \times \mathbb N|).$$

If $U$ maps onto $V$ then $U \times \mathbb N$ maps onto $W=((U \times \mathbb N) \times \{0\}) \cup (V \times \{1\})$ and of course $U \times \mathbb N$ injects into $W$.  Thus $U \times \mathbb N$ is
the same size as $W$, and since $V$ obviously injects into $W$, $V$ injects into $U \times \mathbb N$.  Now consider $Y = X \times \mathbb N$.  We have $Y \sim Y \times \mathbb N$.  A well-ordering of a partition of $Y$ gives a map from a well-ordered set (the partition itself) onto $\iota``Y$, which is in bijection with $\iota``Y \times \mathbb N$, and so an injection from the same well-ordered set into $\iota``Y$, from which we derive a well-ordering of a subset of $Y$ of a corresponding length mod a shift of type.  This justifies $\aleph(X \times \mathbb N)\geq \aleph^*(X \times \mathbb N)$, and the converse inequality is obvious.

We refine the strategy of the first argument to refute the Weak Partition Principle in NF.

Let $X$ be a set.  Let $Y$ be a well-orderable subset of ${\cal P}(X)$.  Let $Z_\alpha$ be an indexing of $Y$. 

For each $x,y \in X$ such that some set in $Y$ contains one but not the other, let $R_{x,y}$ be the first $Z_\alpha$ such that $\neg(x \in Z_\alpha \leftrightarrow y \in Z_\alpha)$, the first set which treats $x$ and $y$ distinctly.  Note that $R_{x,y}=R_{y,x}$.  Let $W$ be a well-ordering of the $R_{x,y}$'s (determined by restricting $Z$):  $$W = \{(A,B) \in Z:(\exists xyuv:A=R_{x,y} \wedge B=R_{u,v})\}.$$

There is a map from $\iota``X^2$ onto the set of $R_{x,y}$'s, sending $\{(u,v)\}$ to $W_0$ if $R_{u,v}$ is not defined, and otherwise sending $\{(u,v)\}$ to $R_{u,v}$.  $W$ is seen to be the order type of a partition $P$ of $X^2$:  for each $\alpha>1$, $P_\alpha$ is $\{(x,y):R_{x,y} = Z_\alpha\}$; $P_0$ contains all other elements of $X^2$.   $T^2(|{\tt dom}(W)|)$ is the cardinality of the set
of order types of initial segments of this partition of $X^2$, so we see that $T^2(|{\tt dom}(W)|)\leq T^3(\aleph^*(|X^2|))$, so $|{\tt dom}(W)| \leq T(\aleph^*(|X|^2)$.

Consider for
each $x \in X$, the set $\chi(\iota^2(x))$ of all $W_\alpha$ such that $x \in W_\alpha$.  Notice that if $x$ and $y$ belong to different
sets in $Y$, $\chi(\iota^2(x))$ and $\chi(\iota^2(y))$ will be different, because the first set which contains one and does not contain the other
will be $R_{x,y}$, a $W_\alpha$.   

Now for each set $A \in Y$, define $\xi(\iota^2(A))$ as the set of $\chi(\iota^2(x))$ such that
$x \in A$.  $\xi$ is clearly an injection:  any element which separates two elements of $Y$ has an analogue via $\chi$ and will separate the analogues of these sets via $\xi$.  

The range of $\xi$ inherits a well-ordering $Q$ from $Y$.  This is a well-ordering of sets of sets of elements of
${\tt dom}(W)$.   $Q = \{(\xi(\iota^2(Z_\alpha)),\xi(\iota^2(Z_\beta))):\alpha \leq \beta\}$ has order type ${\tt ot}(Q) = T^2({\tt ot}(Z))$.

This tells us that any order type of a well-ordering of subsets of $X$ has image under $T^2$ which is the order type of a well-ordering of sets of sets of elements of a well-ordered partition $P$ of $X^2$.

So we have $\{{\tt ot}(W^{\iota^2}):{\tt dom}(W) \subseteq {\cal P}(X)\} \subseteq \{{\tt ot}(Q): {\tt dom}(Q) \subseteq {\cal P}^2(P)\}$

The left side of this inequality has cardinality $T^4(\aleph(T(\exp(|X|))))$ = $\aleph(\exp(T^5(|X|)))$.

The right side of the inequality has cardinality $T^2(\aleph(T(\exp(T(\exp(|P|))))))$ = $\aleph(\exp^2(T^4(|P|)))$.     $T^2(|P|)$ is the cardinality of the set of order types of initial segments of the given well-ordering of $P$ which is less than $T^3(\aleph^*(|X|^2))$, the cardinality of the set of all order types of well-orderings of partitions of $X^2$, so we have $|P|<T(\aleph^*(|X|^2))$, and further $$\aleph(\exp^2(T^4(|P|)))\leq \aleph(\exp^2(T^5(\aleph^*(|X|^2)))) = \aleph(\exp^2(\aleph^*((T^5(|X|)^2)))).$$

When the stratification dust settles, we find that $\aleph(\exp(|X|))$ (the cardinal bound associated with the order type of the well ordering on $Y$, projected down to cardinals at the level of $|X|$) is less than or equal to $\aleph(\exp^2(\aleph^*(|X|^2)))$, the cardinal bound associated with the order type of a well-ordered collection of sets of sets of elements of a well-ordered partition of $X^2$, again all projected down to cardinals at the level of $|X|$.  We have shown all the supporting calculations above:  this is the kind of discipline which would be required in the imaginary community of mathematical researchers who take NF as the ultimate foundation.
$\aleph(\exp(|X|)) \leq \aleph(\exp^2(\aleph^*(|X|^2)))$ is the inequality discovered and proved by ChatGPT Sol.

This then tells us that for any large enough concrete $n$ that all computations can be carried out, $\aleph(\exp(T^n(|V|))) \leq \aleph(\exp^2(\aleph^*(|T^n(|V|)^2)))$.



But then if we assume the Weak Partition Principle $$\aleph(T^{n-1}(|V|)) = \aleph(\exp(T^n(|V|))) \leq \aleph(\exp^2(\aleph^*(|T^n(|V|^2)))) = $$ $$\aleph(\exp^2(\aleph^*(|T^n(|V|)))) = \aleph(\exp^2(\aleph(T^n(|V|))))$$
The last equation uses the fact that $\iota^n``V \sim \iota^n``V \times \mathbb N$ in combination with WPP.

We have then have a contradiction by Henson's lemma  with $\alpha = \aleph(T^{n-1}(|V|))$ and $F(\alpha)$ defined as $\aleph(\exp^2(\alpha))$.   

\section{The proof of Infinity}

The aim here is to give a short version of the proof that Infinity holds in NF, which is ultimately due to Specker.
There are two purposes.  One is to disentangle the proof of Infinity from the more general proofs of failure of Choice.
The other is to fix an issue in NF pedagogy.  The properties of the ordered pair commonly used in NF depend on Infinity holding.  But the proof as usually presented seems to require functions and relations, and so to require a pair.  To avoid using the Kuratowski pair to develop functions and relations and then redefining all these basic notions, we show that the proof actually can be conducted with no pairs at all.

We will develop all concepts used in the proof from the ground up.

This proof looks quite different from the proofs above because of the minimal toolkit used.  We do not even use the order on the natural numbers in this proof (which was surprising to us:  we expected to have to use order).  We will make some comparisons between this argument and the arguments given above at the end of the section.:  we argue that in spite of appearances, it is an argument of the same general sort as the arguments in the proofs above.

\begin{description}

\item[empty set:]  The empty set $\emptyset$ is defined as $\{x\mid x=x\}$.

\item[singleton:]   The singleton $\{x\}$ is defined as $\{y\mid y=x\}$.  We also use the notation $\iota(x)$ for the singleton of $x$, which has the convenient property that it can be iterated.

\item[boolean operations:]  We define $A \cup B$ as $\{x \in A \vee x \in B\}$.  We define $A \cap B$ as $\{x \mid x\in A \wedge x \in B\}$.  We define $A - B$ as $\{x \in A \wedge x \not\in B\}$.  We define $V$ as
$\{x\mid x=x\}$ and $A^c$ as $V-A$.

\item[zero:]  The first natural number 0 is defined as $\{\emptyset\}$.  Notice that, informally, this is the set of all sets with zero elements.

\item[enhancement of set notation:]  We are taking stratified set abstracts as understood.  We define $\{x \in A \mid \phi\}$ as $\{x\mid x \in A \wedge \phi\}$, as usual.  For any term construction $\{F(x_1,\ldots,x_n)$, we define
$\{F(x_1,\ldots,x_n)\mid\phi\}$ as $\{x\mid (\exists x_1\ldots x_n:x=F(x_1,\ldots,x_n) \wedge \phi)\}$.

\item[successor:]  We define $\sigma(A)$ for any set $A$ as $\{a \cup \{x\}\mid a \in A \wedge x \not\in A\}$.

\item[numerals:]  We define 1 as $\sigma(0)$, 2 as $\sigma(1)$ and so forth.  Notice that for any concretely given natural number $n$, the number $n$ thus defined will contain exactly the sets with $n$ elements.

\item[inductiveness:]  The collection $\cal I$ of inductive sets is defined as $$\{I\mid 0 \in I \wedge (\forall A:A \in I \rightarrow \sigma(A) \in I)\}.$$

\item[intersections and unions:]  We define $\bigcup A$ as $\{x\mid (\exists a \in A:x \in a)\}$.   We define $\bigcap A$ as $\{x\mid (\forall a\in A:x \in a)\}$.

\item[natural numbers and finiteness:]  We define $\mathbb N$, the set of natural numbers, as $\bigcap \cal I$.  We define $\mathbb F$, the set of finite sets, as $\bigcup \mathbb N$.

\item[status of the Peano axioms:]  We have defined notions of zero and successor.   $\sigma(A) \neq 0$ is evident, as each element of a set $\sigma(A)$ must be nonempty, and the only element of 0 is empty.
Mathematical induction is evident from the form of the definition of $\mathbb N$.  The status of Peano's fourth axiom remains to be seen.

\item[Lemma:]  Each natural number is either 0 or a successor.  This is obvious:  the set consisting of 0 and all successors of sets is trivially inductive.

\item[Predecessor Lemma:]  For any natural number $n$, for any $A \in \sigma(n)$, for any $x \in A$, $A - \{x\} \in n$. 

\item[Proof:]  For $n=0$, this reduces to $\{x\}-\{x\} = \emptyset \in 0$. 

Fix $k$.  Suppose that for any $B \in \sigma(k)$, for any $y \not\in B$, $B -\{y\} \in k$.  Let $A \in \sigma(\sigma(k))$ and $x \in A$.
$A=C \cup \{z\}$ for some $C \in \sigma(k)$ and $z \not\in C$.  If $x=z$, we have $A - \{x\} = C \in \sigma(k)$.  If $x \neq z$, we have $A - \{x\} = (C \cup \{z\}) - \{x\} = (C - \{x\}) \cup \{z\}$.  $C - \{x\}$ is in $k$ by ind hyp, ao 
$A - \{x\} = (C - \{x\}) \cup \{z\} \in \sigma(k)$.  This completes the proof by induction of the Lemma.

\item[Disjointness Lemma:]  Distinct natural numbers are disjoint.

\item[Proof:]  0 is disjoint from other natural number because the only element of 0 is empty and any element of a successor is nonempty.

Suppose $k$ is disjoint from all other natural numbers.  $\sigma(k)$ is disjoint from 0 for reasons already stated.  Suppose that $\sigma(k)$ and $\sigma(m)$ both contain a set $A$.  This set will be nonempty with an element $x$.  $A-\{x\}$ will belong to both $k$ and $m$ by the Predecessor Lemma, so $k=m$ by ind hyp, so $\sigma(k)=\sigma(m)$.

\item[Peano's fourth axiom for nonempty natural numbers:]  If $\sigma(m)=\sigma(n)$ are nonempty, then $m=n$.

\item[Proof:]  Let $\sigma(m)=\sigma(n)$ and $A \in \sigma(m)=\sigma(n)$.  Then there is $x \in A$ and $A-\{x\}$ belongs to both $m$ and $n$ by the Predecessor Lemma, so $m=n$ by the Disjointness Lemma.

\item[Doubling operation:]  For each natural number $n$, for any element $A$ of $n$ and any $x \not\in \bigcup(A)$, there is a natural number $d(n)$ not depending on choice of $A$ or $x$ such
that $A \cup \{a\mid x \in a \wedge a-\{x\} \in A\} \in d(n)$.  Further, $d$ is injective where defined.

This is evident for $n=0$, with $d(n)=0$, since $A=\emptyset$, $x$ can be anything, and $A \cup \{a\mid x \in a \wedge a-\{x\} \in A\} =0 \in d(0)$ with these assignments of values.

Suppose that for any element $B$ of $k$ and $y \not\in \bigcup B$, we have $B \cup \{a\mid y \in a \wedge a-\{y\} \in B\} \in d(k)$.  Let $A \in \sigma(k)$ and $x \not\in \bigcup A$.  $A = C \cup \{z\}$ for
some $C \in k$ and $z \not\in C$.  $A \cup \{a\mid x \in a \wedge a-\{x\} \in A\} = C \cup \{a\mid x \in a \wedge a-\{x\} \in C \}\cup \{z\} \cup \{z \cup \{x\}\} \in \sigma(\sigma(d(k)))$, so we have established the theorem
and the identity $d(\sigma(k))= \sigma(\sigma(d(k)))$.

Note that $d(k)$ may fail to be defined if every element of $k$ has set union the entire universe.

Notice also that $d(\sigma(k))=d(\sigma(l))$ implies $d(k)=d(l)$ by the limited Peano's fourth axiom.  This implies that $d$ is injective.  Certainly if $d(n)=d(0)$ it follows that $n=0$.  Suppose for a fixed $k$
that if $d(n)=d(k)$ we can deduce $n=k$.  Then if $d(\sigma(k))=d(m)$ we can deduce $m \neq 0$, so $m=\sigma(p)$, so $d(\sigma(k))=d(\sigma(p))$ so $d(k)=d(p)$, so $k=p$, so $\sigma(k)=\sigma(p)=m$.

\item[Images:]  For any term construction $F(x)$ which is stratified, define $F``A$ as $\{F(x):x \in A\}$.

\item[Exponential operation:]  For each natural number $n$ with element $A$, there is a unique natural number $\exp^*(n)$, not depending on $A$, such that ${\cal P}(A)\in \exp^*(n)$.  Further, $\exp^*$ is injective.

\item[Proof:]  For $n=0$, there is only one element $\emptyset$ of the natural number 0, whose power set is $\{\emptyset\}$ which belongs to 1, so $\exp^*(0)=1$.

Suppose for a fixed natural number $k$ that we have determined a unique $\exp^*(k)$ as above.  Let $A = B \cup \{x\} \in \sigma(k)$ with $B \in k$.   Then ${\cal P}(B) \in \exp^*(k)$
and ${\cal P}(A) = {\cal P}(B) \cup \{a\mid x \in a \wedge a-\{x\} \in {\cal P}(B)\} \in d(\exp^*(k))$, so we have proved the existence of $\exp^*(\sigma(k))$ and the identity $\exp^*(\sigma(k))=d(\exp^*(k))$.

Please note that $\exp^*(k)$ is one type higher than $k$ for purposes of stratification.

Now notice that $\exp^*(\sigma(m))=\exp^*(\sigma(n))$ implies $d(\exp^*(m))=d(\exp^*(m))$ which implies $\exp^*(m)=\exp^*(n)$ by injectivity of $d$.  From this it follows that $\exp^*$ is injective.
Certainly $\exp^*(n)=\exp^*(0)=1$ implies $n=0$.  Suppose that for a fixed $k$, $\exp^*(n)=\exp^*(k)$ implies $n=k$ for all $n$.  Then $\exp^*(\sigma(k)) =\exp^*(m)$ implies first that $m \neq 0$, so we can take $m=\sigma(p)$, so $\exp^*(\sigma(k))=\exp^*(\sigma(p))$, so $\exp^*(k)=\exp^*(p)$, so $k=p$, so $\sigma(k)=\sigma(p)=m$.

\item[Type raising:]  For each nonempty natural number $n$, there is a unique natural number $T(n)$ such that for all $A$, $\iota``A \in T(n)\leftrightarrow A \in n$.

\item[Proof:]  $T(0)=0$ because $\iota``\emptyset = \emptyset$.

If $T(k)$ is defined and $A = B \cup \{x\} \in \sigma(k)$, with $B \in k$ and $x \not\in B$, then $\iota``A = \iota``B \cup \{\{x\}\} \in \sigma(T(k))$.  
If $\iota``C \in \sigma(T(k))$ then $\iota``C = \iota``D \cup \{\{z\}\}$ for some $\{z\} \not\in \iota``D$, so $\iota``D \in T(k)$, so $D \in k$, so $C \in \sigma(k)$.
So $T(\sigma(k))$ exists and is equal to $\sigma(T(k))$.

This does not support a proof that $T(k)=k$ for all $k$, because we cannot form the set of all $k$ satisfying this unstratified condition.

\item[Bijection lemma:]  If $F(x)$ is a term construction of the same type as $x$ with $F(x)=F(y) \rightarrow x=y$ in all cases, then for any natural number $n$ and set $A$ on which the term construction
is defined for each element. $A \in n \leftrightarrow F``A \in n$.

\item[Proof:]  For $n=0$, note that $F``\emptyset = \emptyset$.

Suppose the result shown for $k$.  Let $A = B \cup \{x\} \in \sigma(k)$ with $B \in k$ and $x \not\in B$.  We have $F``A = F``B \cup \{F(x)\}$.  $F``B \in k$ by ind hyp, and $F(x) \not\in F``B$ because
$F$ is injective, so $F``A \in \sigma(k)$.

\item[Lemma:]  $T(\exp^*(n)) = \exp^*(T(n)$

\item[Proof:]   Let $A \in n$.  The first number is the number of  singletons of subsets of $A$.  The second number  is the number of sets of singletons of elements of $A$.
The construction $F(x) = \iota``\bigcup(x)$ on sets of singletons and $x$ otherwise is injective and maps singletons of subsets of $A$ onto subsets of $\iota``A$.  So the two numbers
are the same by the Bijection Lemma.

\item[Exponential function:]  We define $\exp(n)$ for a natural number $n$ as $T^{-1}(\exp^*(n))$.  Notice that this is of the same type as $n$ for stratification purposes.  Also notice that it may be partial:  for $\exp^*(n)$ to be defined, 
an element $A$ of $n$ has to have its power set belonging to the same natural number as a set of singletons.  Notice further that $\exp$ is injective where defined.

\item[Hierarchy of $n$:]  For any natural number $n$, we define $H(n)$ as the smallest set which contains $n$ and is closed under $\exp^{-1}$ (that is, the intersection of all sets which contain $n$ and
for any $m$, contain $m$ if they contain $\exp(m)$).  We aren't presuming here that $\exp^{-1}$ is total.

\item[Subset Lemma:]  Every subset of a finite set is finite.

\item[Proof:]  We prove by induction that for each natural number $n$ and $B \subseteq A \in n$, $B$ is finite.

If $n=0$ and $B \subseteq A \in n$, then $B=\emptyset \in 0$, so $B$ is finite.

Suppose that each element of a fixed natural number $k$ is finite.  Let $A \in \sigma(k)$, $x \in A$, and $B \subseteq A$.  $B-\{x\} \subseteq A-\{x\}\in k$ is finite, so $B = (B-\{x\})\cup \{x\}$ is finite (whether $x \in B$ or not).

\item[Corollary:]  If $V$ is finite, then every set is finite.

\item[Proper Subset Lemma:]  For each natural number $n$, no element of $n$ has a proper subset belonging to $n$.

\item[Proof:]   The sole element $\emptyset$ of 0 has no proper subsets.


Suppose for a fixed natural number $k$ that no element of $k$ has a proper subset in $k$.  Suppose that $A$ is a proper subset of $B$ and both belong to
$\sigma(k)$.  Let $x \in A$.  Then $A-\{x\}$ is a  proper subset of $B-\{x\}$ and both belong to $k$, which is a contradiction,  completing the proof of the Lemma.


\item[Corollary:]  If $V\in n \in \mathbb N$ then $n=\{V\}$ and $\sigma(n)=\emptyset$.  Note that it is straightforward to see that $\emptyset$ can be a natural number only if its predecessor is $\{V\}$, so in 
fact $\emptyset \in \mathbb N$ is a (rather mysterious) equivalent of the negation of Infinity.

\item[Definition:]  A natural number $n$ is {\em even\/} if it is $d(m)$ for some $m$;  otherwise it is {\em odd\/}.   If two natural numbers are both even or both odd, we say they have the same parity;  otherwise we say the two natural numbers have different parities.

\item[Parity Lemma:]  The successor of an even number is odd if nonempty;  The successor of an odd number is even if nonempty.  $n$ is even iff $T(n)$ is even.

\item[Proof:] 0 is even and its successor 1 is odd:  it should be clear that 1 is odd because for any $A$ and any $x \not\in \bigcup A$, the set $A \cup \{a:a-\{x\} \in A \wedge x\in a\}$ has at least two distinct elements if it has any elements.
1 is odd and 2 is even ($2=d(1)$).

Suppose for a fixed $k$ that $k$ is even iff $\sigma(k)$ is odd.

We argue that for any $n$, $\sigma(\sigma(n))$ is even iff $n$ is even (if $\sigma(\sigma(n))$ is nonempty).  If $n$ is even then $n=d(m)$ and $d(\sigma(m))=\sigma(\sigma(d(m))=\sigma(\sigma(n))$ so $\sigma(\sigma(n))$ is even.

There is a boundary case here:  what if $d(\sigma(m))$ is undefined?  This means that every element $A$ of $\sigma(m)$ has $\bigcup A=V$.  Since $d(m)$ is defined by hypothesis, there is $B \in m$ with
$y \not\in \bigcup B$.  It follows that every element of $V-B$ must contain $y$, and we see that $d(m)$ in fact contains the universe:  $V-B$ is precisely the set of things obtained by adding $y$ to elements of $B$, and the claim follows from the definition of $d$.  So $V \in n = d(m)$ in the boundary case and $\sigma(\sigma(n)$ is empty.

  If $\sigma(\sigma(n))$ is even
then it is $d(m)$ for some $m$, and this $m$ can be written $\sigma(p)$, because certainly $m$ is not zero.  $\sigma(\sigma(n)=d(m)=d(\sigma(p))=\sigma(\sigma(d(p)))$ whence by modified Peano's fourth $d(p)=n$, so $n$ is even.

Now for $k$ above we have $\sigma(k)$ is even iff $\sigma(\sigma(k))$ is odd (and nonempty).  This establishes the result by induction.

0 is even iff $T(0)$ is even as both are even.

If $k$ is even iff $T(k)$ is even, 
it follows from the first part of the lemma that $\sigma(k)$ is even iff $\sigma(T(k))=T(\sigma(k))$ is even.

\item[Theorem (Infinity):]  $V$ is not finite.

Suppose to the contrary that $V$ is finite and belongs to a natural number $N$.

Observe that $\exp^*(N)=N$ (the universe is its own power set), whence $\exp^*(T(N))=T(\exp^*(N)=T(N)$ and $\exp(T(N))=T^{-1}(\exp^*(T(N))=T^{-1}(T(N))=N$.

Note that $T(N) \neq N$ because $T(N)$ contains the proper subset $\iota``V$ of the sole element $V$ of $N$.  Further, $N$ is not  $T(n)$ for any natural number $n$ for the same reason:  $\iota``A$ is a proper subset of $V$ for $A$ in any natural number $n$.  Further note that $\exp(N)$ is undefined, because $\exp^*(N)=N$ is not an image under $T$.

The set $H(T(N))$ consists of $T(N)$ and all of its inverse images under $\exp$.   

The set $H(N)$ is precisely $H(T(N)) \cup \{N\}$.   Assume $n \in H(N)-\{N\}$.  Suppose $A$ contains $T(N)$ and is closed under $\exp^{-1}$.  Then $A \cup \{N\}$ contains $N$ and is closed under $\exp^{-1}$,
so it contains $n$, so $n \in A$, so $n \in H(T(N))$.   Suppose $n \in H(T(N))$:  then it certainly belongs to any $B$ containing $N$ and closed under $\exp^{-1}$, so it is in $H(N)$.  So we see that
$H(T(N)$ is precisely $H(N)-\{N\}$. 

$H(T(N))$ belongs to some natural number $M$ (on the assumption that $V$ is finite)
and $H(N)$ is seen to belong to $\sigma(M)$.

But, also, $H(T(N)) = T``H(N)$.    $T(N) \in T``H(N)$ and $T``H(N)$ is closed under $\exp^{-1}$ so all elements of $H(T(N))$ are in $T``H(N)$.  The image of $N$ under $T$ is in $H(T(N))$ and if $T(k)$ is in $H(T(N))$, it follows
that $T(\exp^{-1}(k))=\exp^{-1}(T(k))$ is in $H(T(N))$, so each element of $T``H(N)$ is in $H(T(N)$.

$T``H(N)$ belongs to $T(\sigma(M))$ because $T``H(N)$ is easily seen to be in bijection with $\iota``H(N)$:  the operation taking $T(n)$ to $\{n\}$ for all natural numbers $n$ and fixing all other objects is type level and the Bijection Lemma shows the needed result.

But this contradicts the Parity Lemma, because $M$ and $\sigma(M)$ have different parities, while $\sigma(M)$ and $T(\sigma(M))$ have the same parity.

The proof of the Infinity Theorem is complete.
\end{description}

We need to do a little extra work to show that this proof actually belongs to the same family as the proofs above.  We will articulate a result clearly similar to Henson's lemma of which the closing stages of the argument above can be seen to be an instance.

In the argument that follows to the end of this section, we do not necessarily give all details.  In particular, the reader may need to supply, if desired, verification that sets which we claim are orbits are in fact orbits.  This is not a problem for us, because our aim is not to prove Infinity (we already did this) but to illustrate the relationship between its proof and the proofs above.

\begin{description}

\item[Definition (order on the natural numbers):]  $m \leq n$ is defined as holding iff each element of $n$ is a superset of some element of $m$.   We define $<, >, \geq$ in the obvious ways.

\item[Theorem:]  $\leq$ is a well-ordering.

\item[Proof:]  That $\leq$ is reflexive and transitive is evident from the definition.

We need to show that $\leq$ is antisymmetric.  Suppose that $m \leq n$ and $n \leq m$.  Then each element $A$ of $m$
is a superset of an element $B$ of $n$ which must in turn be a superset of an element $C$ of $m$.  But by the proper subset
lemma, $A=B=C$.  The argument is symmetric, so $m$ and $n$ have the same elements, and $m=n$.

We need to show that $\leq$ is total.  We prove by induction on $n$ that $m \leq n$ or $n\leq m$.  For $n=0$ this is evident because $0 \leq m$ is obvious.  Suppose we have shown this for $k$.  We have $m \leq k$ or $k \leq m$.  Clearly
$k \leq \sigma(k)$, so if $m \leq k$ we have $m \leq \sigma(k)$.  Suppose $k \leq m$.  If $k=m$ then $m \leq \sigma(k)$.
If $k \neq m$, then each element of $m$ is a proper superset of an element of $k$, so includes an element of $\sigma(k)$ as a subset, so $\sigma(k) \leq m$.

We need to show that each nonempty set of natural numbers has a $\leq$-minimal element.  Suppose $A$ is a set of natural
numbers with no $\leq$-minimal element.  We show that $A$ is empty.  What we actually show by induction is that for each
$n$, for all $m \leq n$, $m \not\in A$.  $0 \not\in A$, because if it did belong to $A$ it would clearly be $\leq$-least.
Thus for all $m \leq 0$, $m\not\in A$.  Now suppose that for all $m \leq k$, $m \not\in A$.   Suppose $\sigma(k) \in A$.
By hypothesis, $\sigma(k)$ is not $\leq$-least, so we have some $p$ such that $k < p < \sigma(k)$.  But clearly $p<\sigma(k)$ implies $p \leq k$:  any element of $\sigma(k)$ properly including an element of $p$ implies that there is an element of $k$ including the element of $p$, so $\sigma(k)>p$ implies $k \geq p$.

\item[Lemma:]  Let $m$ and $n$ be natural numbers.  If any element of $n$ is a superset of an element of $m$, then $m \leq n$.

\item[Proof:]  We prove this by induction on $n$.  For $n=0$ it is trivial, since the hypothesis implies $m=0$.
Suppose it true for $k$.  Suppose that an element of $\sigma(k)$ is a superset of an element of $m$.  If $m=\sigma(k)$
we have $m \leq \sigma(k)$ as desired.  Otherwise, the given element of $\sigma(k$ is a superset of an element of $k$ which includes an element of $m$ as a subset, so by ind hyp $m \leq k$, so $m \leq \sigma(k)$ by transitivity.

\item[Lemma:]  For nonempty natural numbers $m$ and $n$, $m \leq n$ iff $T(m) \leq T(n)$.

\item[Proof:]  It is sufficient to observe that $A \leq B$ iff $\iota``A \subseteq \iota``B$.

\item[Definition (orbit, iteration):]  Let $F(x)$ be a term construction of the same type as $x$.   We say that $O$ is an $F$-orbit from
$x$ to $y$ of length $n$ iff 

\begin{itemize}

\item $x \in O$, $y \in O$, 

\item for each $z \in O$ other than $y$, $F(z) \in O$, 

\item and for each $z \in O$ other than $x$ there is exactly one $w\in O$ such that $F(w)=z$,

\item no proper subset of $O$ has the properties above,

\item and $O \in T^2(n)$.
\end{itemize}

\item[Lemma:]  For each $x$ and $n$, there is no more than one $y$ such that there is an $F$-orbit from $x$ to $y$ of length $n$.

\item [Proof:]  We prove this by induction on $n$.  For $n=0$, it is clear that $\{x\}$ is an $F$-orbit from $x$ to $x$ of length 0,
and moreover $y=x$ is the only $y$ for which there is an $F$-orbit from $x$ to $y$ of length 0.

Suppose that the result has been shown for $n=k$.   There is no more than one $y'$ such that there is an $F$-orbit $O'$ from
$x$ to $y'$ of length $k$.  Suppose there is such a $y'$.  Set $y=F(y')$ and $O = O' \cup \{y\}$.  Then $O$ is an $F$-orbit from $x$
 to $y$ of length $\sigma(k)$ unless $y \in O$.
 
 Suppose that $O$ is an $F$-orbit from $x$ to $y$ of length $\sigma(k)$.  Then there is $z \in O$ such that $F(z)=y$,
 and $O - \{y\}$ is an $F$-orbit from $x$ to $z$ of length $k$, so $z=y'$ and $y=F(y')$.
 
 We have established that at most one such $y$ can exist.
 
 \item [Definition:]  We define $F^n(x)$ as the unique $y$ such that there is an $F$-orbit from $x$ to $y$ of length $n$.
 
 \item[Baby Henson lemma:]  Let $F(x)$ be a term construction of the same type as $x$.  Let $F(T(n))=T(F(n))$ hold for all natural numbers, and let $F$ be nondecreasing (if $m \leq n$ then $F(m) \leq F(n)$ if both are defined). Let $M$ be such that $T(M) < M$.  Then it cannot be the case that $F(T(M))=M$.
 
 \item[Proof:]   Suppose to the contrary that $F(T(M))=M$.
 
 Let $m$ be smallest such that there is an $F$-orbit $O$ from $m$ to $M$ of some natural number length $k$.
 Notice that there will be such an $m$ because $T(M)$ satisfies the condition.  So we have $F^k(m)=M$.
 
 Observe that $T``O$ is an $F$-orbit from $T(m)$ to $T(M)$ of length $T(k)$ [verification that it is an orbit relies on $F$ commuting with $T$; note that $T``O$ belongs to $T(\sigma(k))$ because there is a bijective type level relation between $\iota(x)$ and $T(x)$].  Further, $T``O \cup \{M\}$ is an $F$-orbit from
 $T(m)$ to $M$ of length $\sigma(T(k))$.  This implies $T(m) \geq m$ because $m$ is minimal.  Further, if $T(m) > m$
 it follows that $T^{-1}(m)$ exists and is less than $m$, and $T^{-1}``(O \setminus \{M\})$ would be an $F$-orbit from
 $T^{-1}(m)$ to $M$, contradicting minimality of $m$.  Thus $T(m)=m$.
 
 Because $T(m)=m$, we see that $T``O \cup \{M\}$ and $O$ must be the same orbit.  Thus $\sigma(T(k))=k$, which is impossible for reasons of parity.
 
 \item[Theorem:]  $V$ is not finite.
 
 \item[Proof:]  It is straightforward to see that $F(x)=\exp(x)$, $M=N$ satisfy the conditions of the baby Henson's lemma.
 
 \item[Observation:]  The proof given first is actually a streamlined version of the same argument.  For in the first proof of Infinity the set $H(N)$ is actually
 the $\exp$-orbit from the minimal $m$ for which such an $\exp$-orbit exists to $N$.  In the argument as given first, it is not necessary to say anything about the minimal value $m$ (or about order at all).  This is interesting but does not change the genealogy of the argument.
 
 
 
 

\end{description}


\section{Specker's original version of the proof of the negation of Choice compared with the current proofs}

We give Specker's original argument.  He defined the set of cardinals which we here will call $SM$ (for {\em Speckermenge\/}) which is the set of cardinals with finitely many iterated images under $\exp$.  Assume AC.  Let $\mu$ be the smallest of the {\em Speckermenge\/}, and let $\exp^n(\mu)$ be the highest indexed iterated image of $\mu$ under $T$.  Note that $\exp^n(\mu) \geq T(|V|)$, since it has no image under exp.  $T(\mu) \in SM$, because $\exp^{T(n)}{T(\mu)} = T(\exp^n(\mu)) \geq T^2(|V|)$, which has either 1 or 2 further iterated images under exp.
$\exp^{Tn+1}(\mu)$ exists and is $>T(|V|)$ (otherwise $\exp^{n+1}(\mu)$ would be defined contrary to assumption,
and $\exp^{Tn+2}(\mu)$ either fails to exist or is $|V|$ exactly.  Further $T^{-1}(\mu)$ exists, because $\mu<T^2)(|V|)$,
and $T^{-1}(\mu) \in SM$.  We have $\exp^n(\mu)>T(|V|$, because its image under exp is undefined, so
$\exp^{n-2}(\mu)>T^3(|V|)$ (whose double image under exp is $T(|V|)$.  This allows us to see that
$\exp^{T^{-1}(n)-2}(T^{-1}(\mu))>T^2(|V|)$, which means that $T^{-1}(\mu)\in SM$ as well:  it can have at most two
further iterated images under exp.  A point which we do not recall being discussed is the fact that $T^{-1}(n)-2$ actually exists:  this is because we are given a collection of $n-2$ cardinals $\exp^i(\mu)$ which are all images under $T$.  Now if $\mu$ is not cantorian, one of $T(\mu)$ and $T^{-1}(m)$ will be less, contradicting minimality of $\mu$ in SM.  Thus $\mu=T(\mu)$, and this means that $n$ is equal either to $T(n)+1$ or $T(n)+2$, which is absurd because the residues of $n$ and $T(n)$ mod 3 are the same.

The proof of the failure of Choice from Henson's lemma is immediate:  let the special function $F(x)$ be defined as $\exp(x)$ and let $\alpha=T^n(|V|)$.  This satisfies the conditions of Hensons lemma and we obtain a contradiction from
$\exp(T^n|V|)=T^{n-1}(|V|)$.

What is strange here is that the involvement of mod 3 in the Specker proof vanishes.  It should be noticed that this is not really a new observation: it is implicit in \cite{henson} though not particularly noted that I recall.  But subsequent workers in NF have not picked up on the fact that Henson turned the proof into a parity argument.

To make the point more explicitly, we give the proof directly.  We find it much preferable to use $\alpha=T(|V|)$ than $\alpha=|V|$ for reasons the reader may divine.

Let $\mu$ be the smallest cardinal such that $\exp^n(\mu)\geq T(|V|)$ for some $n$.  Let $n$ be minimal.
We claim that $T(\mu)$ is the smallest cardinal such that $\exp^m(\mu) \geq T^2(|V|)$ for some $m$.  Suppose
that $\nu < T(\mu)$ has $\exp^m(\mu) \geq T^2(|V|)$ for a given $m$.  Then $\nu = T(\kappa)$ for some $\kappa<\mu$
and we have $\exp^{T^{-1}(m)}(\kappa) \geq T(|V|)$, contrary to minimality of $\mu$.  $T^{-1}(m)$ exists because
$m$ is the length of a finite sequence of cardinals which are images under $T$.   Now we see that
$\exp^{T(n)}(T(\mu))\geq T^2(|V|)$, so $\exp^{T(n)+1}(T(\mu))\geq T(|V|)$, so $T(\mu)\geq \mu$.
Further, $\exp^n(\mu)\geq T(|V|)\geq T^2(|V|)$, so $T(\mu) \leq \mu$.   So $T(\mu)=\mu$.  Now
we have $\exp^{n-1}(\mu)<T|V|$ and $\exp^n(\mu)\geq T(|V|)$.  We have $\exp^{T(n)-1}(T(\mu)) < T^2(|V|)$,
$\exp^{T(n)}(T(\mu)) \geq T^2(|V|)$, and $\exp^{T(n)+1}(T(\mu)) \geq T(|V|)$.  So $n=T(n)$ or $T(n)+1$.
But it cannot be $T(n)$ because $\exp^{T(n)}(T(\mu))=T(\exp^n(\mu))$ and this cardinal cannot be cantorian.
So $n=T(n)+1$, which is impossible for reasons of parity.


\section{Specker trees}




\end{document}